\chapter{Deriving an upper bound on the Zeta zero count using Turing's method}
\label{cha:n_upper_bound}

This chapter derives an upper bound of $N(T)$ when $T$ is a certain type of
Gram point which, combined with a lower bound coming from numerical zero
counting on the critical line, forces $N(T)$ to take a value. 
For each natural $n$ the Gram point is the unique positive real $g_n$ such that 
$\theta(g_n) = n\pi$.
This is a strictly increasing sequence which are often, although not always, 
boundaries of intervals containing exactly one zero of the Riemann zeta 
function on the critical line.
One way to gain an upper bound on $N(T)$ is to find a long enough
sequence of Gram points which sequentially bound at least a single 
zeta zero on the critical line.
"Long enough" here is determined by finding explicit constants bounding 
the average behaviour of the zeta argument function
\[
  \left|\int_{t_1}^{t_2} S(t) dt\right| \leq A + B * \log(t_2/2\pi).
\]
This is the work of Ch.~\ref{cha:lehman_sharp}, which proves
$A = 5$ and $B = 0.128$.
In \cite{Lehman1970}, the full result is $A = 2.3$ and $B = 0.128$.
Further improvements have since been made (such as \cite{TrudgianTuring}).

As a result, the main theorem of this section is as follows.

\begin{theorem}
\label{thm:n_upper_bound}
\lean{riemannZeta.N_le_of_gramPoints}
\leanok
Let $n$ and $k$ be natural numbers such that $g_n>168\pi$ and
that $g_n$, $g_{n + k + 1}$ satisfy
\[
  (5 + 0.128\log(g_{n + k + 1} / 2\pi)) < \frac{g_{n + k + 1} - g_n}{2}.
\]
Suppose, moreover, that for each $0\leq i \leq k$ there is at least one
$t\in[g_{n + i}, g_{n + i + 1})$ such that $\zeta(1/2 + it) = 0$.
Then 
\[
  N(g_n) \leq n + 1.
\]
\end{theorem}
\begin{proof}
\leanok
\uses{lem:le_zerofree_dense, thm:ls_five, thm:N_formula, thm:ga_main}
Let $x(t) := (\theta(t) - \theta(g_n)) / \pi$. 
If $t \in [g_{n + j}, g_{n + j + 1})$, then $\lfloor x(t)\rfloor = j$, 
which follows by strict monotonicity of $\theta$ and the fact that 
$\theta(g_j) = j\pi$. 
By the assumption that each Gram interval contains a zeta zero, we have 
$j \leq N(t) - N(g_n)$ for each $t \in [g_{n + j}, g_{n + j + 1})$.
By the zero-free density lemma (Lemma~\ref{lem:le_zerofree_dense}) and 
Theorem~\ref{thm:N_formula} one obtains 
\[
  S(g_n)(g_{n + k + 1} - g_n) \leq \int_{g_n}^{g_{n + k + 1}} \{x(t)\}dt + \int_{g_n}^{g_{n + k + 1}} S(t)dt
\]
Now, $\{x(t)\}$ is one-periodic, antitonic on the period and convex.
Thus its integral is $\leq (g_{n + k + 1} - g_n) / 2$.
By the Lehman-style bound in Theorem~\ref{thm:ls_five}, dividing through by $g_{n + k + 1} - g_n$ we then obtain
\[
  S(g_n) \leq \frac{1}{2} + \frac{5 + 0.128\log(g_{n + k + 1} / 2\pi)}{g_{n + k + 1} - g_n}
\]
Finally, by hypothesis the second term is $< 1/2$, so $S(g_n) < 1$.
Note that $N(g_n) = n + 1 + S(g_n)$ by Theorem~\ref{thm:N_formula}, which 
forces $S(g_n)$ to be an integer and therefore $S(g_n)\leq 0$, 
completing the proof.
\end{proof}


